\documentclass[12pt]{article}
\usepackage[utf8]{inputenc}
\usepackage[T1]{fontenc}
\usepackage{hyperref}
\title{Varli Language Resources for K-12 Curriculum Development}
\author{Compiled by Claude Sonnet 5 (claude-sonnet-5) for Andr\'as Kornai}
\date{}
\begin{document}
\maketitle

\section{Introduction}
Varli (Warli, Varali; ISO 639-3 \texttt{vav}, Glottocode varl1238) is an Indo-Aryan language of the Warli people in Maharashtra (Palghar/Thane area) and neighbouring regions, with about 390,000 speakers in the 2011 census according to Wikipedia. It is usually classed with Marathi, though some scholars have classed it with Konkani or Bhil. It is traditionally unwritten, and Devanagari or Gujarati script is used when written.

\section{Primary Educational Resources}
No primers or textbooks were found. A Rural India Online article reports a primary-school teacher in Mokhada who says Warli is the most widely spoken language there and that teachers bridge English to Warli through Marathi: \url{https://ruralindiaonline.org/ur/articles/palghar-school-teacher-pushes-for-warli}

\section{Linguistic and Reference Materials}
\begin{itemize}
\item Wikipedia overview of Varli: \url{https://en.wikipedia.org/wiki/Varli_language}
\item \emph{Warli of Thana} (Survey of Marathi Dialects VII), an older dialect survey, listed by a book seller (no dictionary found): \url{https://www.exoticindia.com/book/details/warli-of-thana-survey-of-marathi-dialects-vii-old-and-rare-book-hcb879}
\item Outlook India article on a lexicographical exercise collecting Warli terms for climate and environment in Saphale, Palghar (work in progress, not a published dictionary): \url{https://www.outlookindia.com/national/what-do-you-call-climate-change-in-warli-language}
\end{itemize}

\section{Religious and Community Resources}
\begin{itemize}
\item Wikipedia on the Warli people: \url{https://en.wikipedia.org/wiki/Warli_people}
\end{itemize}

\section{Digital and Audio Resources}
None were found in the search results.

\section{Curriculum Development Context}
The language is reported as not taught in schools. No standalone dictionary, grammar or primer was found. The Mokhada teacher report and the Saphale lexicon effort suggest local interest and a starting point for community-led materials, but substantial new documentation would be needed.
\end{document}